\section{Route-A outcome and Route-B lock}
\label{sec:route}

The source lock is complete: the tensor-prime atom set, entropy order,
directed grammar, endpoint roofs, equal edge potentials, $\ell^2$ function
space, transfer $L_s$, chiral pencil $\cB_s$, and the Fredholm/$\det_3$
conventions are fixed before interpretation.  The construction is assigned
candidate identifier SD-C09.

\subsection{Route-A layers}

\begin{center}
\small
\begin{tabularx}{0.99\textwidth}{L{0.08\textwidth}L{0.25\textwidth}X}
\toprule
Layer & Verdict & Evidence \\
\midrule
A0 & analytic arithmetic origin & tensor indecomposables generate primes, entropy orders them and gives $\log p$, without prime or zero tables \\
A1 & analytic pass & only atom loops are periodic; their repetitions give the exact $p^r$ ledger, while successor edges are acyclic \\
A2 & analytic determinant & $\det(I-zL_s)=\prod_p(1-zp^{-s})$ on $\RePart s>1$; signed cancellations are unnecessary \\
A3 & partial analytic structure & $\det_3(I-z\cB_s)$ is holomorphic on $1/3<\RePart s<2/3$, symmetric under $s\mapsto1-s$, and moves on the critical axis, but it is not the Euler determinant continued \\
A4 & formal hint & a canonical self-adjoint critical-axis pencil exists; no fixed generator, quantization, domain theorem, or target counting law is supplied \\
\bottomrule
\end{tabularx}
\end{center}

The tuple is
\begin{center}
\small
\begin{tabular}{c}
\texttt{(A0\_ANALYTIC\_ARITHMETIC\_ORIGIN, A1\_PASS\_ANALYTIC,}\\
\texttt{A2\_ANALYTIC\_DETERMINANT, A3\_PARTIAL\_ANALYTIC\_STRUCTURE,}\\
\texttt{A4\_FORMAL\_HINT).}
\end{tabular}
\end{center}
with overall verdict
\[
 \boxed{\texttt{ROUTE\_A\_ANALYTIC\_CANDIDATE}.}
\]
The stage status is
\[
 \boxed{\textsf{GO A3 CHIRAL MOTION}
 \;/\;
 \textsf{STOP UNIFIED DIVISOR}.}
\]

\subsection{Adversarial boundary}

The argument does not certify RH-like behavior for controls.  Arbitrary DAG
orders and phases, matched nonprime masses, and randomized positive diagonal
masses all retain the all-order triangular determinant mechanism; all 24
frozen random DAGs also show chiral motion.  Only recurrent reverse edges
break the ledger.  Two-state prefixes have
periodic linearly counted crossings.  Thus ``exact diagonal ledger plus
moving chiral singular values'' is not sufficient for a Riemann divisor.
The arithmetic source remains selective at A0, but the chiral inference
stops at A3/A4.

\subsection{Why Route B is locked}

$\cB_{1/2+it}$ is a compact self-adjoint operator for each $t$, but it is an
operator pencil whose parameter is the proposed spectral height.  It is not
a single self-adjoint operator $H$ with eigenvalue parameter $E$.  No dense
domain or boundary condition problem has been posed, no compact-resolvent
$T\log T$ counting theorem exists, no exact von-Mangoldt trace formula is
derived from this pencil, and no spectral determinant equals completed
$\xi$.  The Euler and chiral determinants remain distinct data types.  Route
B therefore cannot be used to combine their best coordinates, and
\[
 \mathtt{route\_b\_invocation\_allowed=false}.
\]

The next smallest same-family test is a relative-determinant theorem: either
derive a source-internal bridge between the Euler Fredholm determinant and
the chiral $\det_3$ without fitted counterterms, or prove that the deleted
low-order traces obstruct every such bridge.
